\documentclass[11pt]{article}
\usepackage{ap-lang}
\title{Part 3: Reasoning and Organization}
\author{Paul Wang}
\date{November 2026}
\docpart{Part 3 \textperiodcentered\ Study Guide}

\begin{document}
\maketitle

\begin{abstract}
Claims and evidence become an argument only when they are arranged into a line of reasoning. This guide covers the line of reasoning itself (Section 1), the methods of development writers use to build paragraphs and whole texts (Section 2), overall structure and where a writer places the thesis, concessions, and turns (Section 3), and transitions, the words that signal logical relationships (Section 4). Section 5 shows how these ideas appear in writing questions, and Section 6 is a method for the synthesis essay (Question 1). Each section ends with check questions; the key is on the last page.
\end{abstract}

\scoringline{Check questions: \blank\ / \qtotal \qquad Target: 13+ \qquad Terms I mix up: \blank[10em]}

\section{The line of reasoning}

A \textbf{line of reasoning} is the sequence of claims a writer uses to lead the audience to the thesis. Each claim is a step; each step should depend on the one before. Reading questions ask how a paragraph advances the line of reasoning; the essay rubric asks you to build one.

\begin{study}{Reasoning as a chain}
Thesis: \emph{The town should replace its annual fireworks with a drone show.}
Step 1: The fireworks have started two brush fires in five years (the problem). Step 2: A drone show produces the same communal spectacle without sparks (the alternative meets the need). Step 3: It costs less after the second year because the drones are reused (the objection about cost is answered). Step 4: Other towns that switched kept their crowds (the fear that tradition will die is answered).
\end{study}

\begin{keyidea}
Ask of every paragraph: \emph{what does the reader now believe that they did not believe before, and why does the next paragraph need it?} If a paragraph could be moved anywhere without loss, it is not doing reasoning work.
\end{keyidea}

\begin{q}{3}{3.1 Line of reasoning}{Answering the cost objection is one step in the chain toward the thesis.}
In the drone-show example above, Step 3 contributes to the line of reasoning by
\ch{changing the subject to the town's budget}{restating the thesis in different words}{answering a likely objection before the audience can raise it}{providing a personal anecdote}
\end{q}

\section{Methods of development}

A \textbf{method of development} is a pattern a writer uses to develop an idea. Most texts mix several.

\begin{terms}
Narration & Telling a story in time order; makes an issue human and memorable. \\
Description & Sensory detail; makes the reader see or feel a place or situation. \\
Exemplification & Supporting a general claim with examples. \\
Comparison and contrast & Placing two things side by side to show likeness or difference. \\
Definition & Explaining what a term means, often to redefine it in the writer's favor. \\
Cause and effect & Tracing how one thing produces another. \\
Classification & Sorting things into categories. \\
Problem and solution & Describing a problem, then proposing a remedy. \\
\end{terms}

\begin{trap}
Definition as argument. When a writer pauses to define a word (``freedom,'' ``success,'' ``waste''), the definition is rarely neutral; it usually reshapes the debate so that the writer's position follows. Questions often ask about that function.
\end{trap}

\begin{q}{4}{3.2 Methods of development}{Tracing how one thing produces another = cause and effect.}
A paragraph explains how a new highway lowered land prices along the old main street, which closed shops, which emptied the downtown. The paragraph is developed mainly through
\chh{classification}{narration}{definition}{cause and effect}
\end{q}

\begin{q}{4}{3.2 Methods of development}{Redefining a key term reshapes the argument in the writer's favor.}
A writer argues that ``a library is not a building full of books; it is the one public room where no one has to buy anything to stay.'' The writer is using definition mainly to
\ch{describe the library's architecture}{give the dictionary meaning of ``library''}{classify the kinds of libraries}{redefine the library's value in a way that supports the argument}
\end{q}

\section{Structure: the shape of the whole}

\subsection{Where the thesis goes}
\begin{compare}{Structure}{How it works}{When writers choose it}
Deductive (thesis first) & State the claim, then support it. & The audience is open or neutral; clarity matters most. \\
Inductive (thesis late) & Build evidence or a story first; the claim emerges near the end. & The audience is resistant; the writer needs to earn agreement before asking for it. \\
Problem--solution & Present the problem, then the proposal. & Policy arguments. \\
Chronological & Follow events in time order, often with a turn. & Narratives, histories, speeches marking an anniversary. \\
\end{compare}

\subsection{Shifts and turns}
Many passages pivot: from story to argument, from past to present, from concession to rebuttal, from criticism to call to action. Questions ask about these turns (``the shift in paragraph 4 from\ldots to\ldots''). Look for words like \emph{but}, \emph{yet}, \emph{today}, \emph{so}, and for changes in pronouns (\emph{I} to \emph{we}) or tense.

\subsection{Openings and closings}
Openings may use an anecdote, a striking statistic, a question, or a scene to establish exigence. Closings may return to the opening image (a \textbf{frame}), broaden the claim, issue a call to action, or end on an unresolved question.

\begin{q}{2}{3.3 Structure}{Delaying the thesis suits an audience that must be won over first.}
A writer addressing an audience opposed to her view presents three stories of affected families and states her proposal only in the final paragraph. This inductive structure is most likely chosen because
\ch{inductive structures are required in formal writing}{the audience must be drawn toward agreement before hearing the request}{the writer is unsure of her own position}{the audience already agrees with her}
\end{q}

\begin{q}{1}{3.3 Structure}{Returning to the opening image frames the essay and signals closure.}
An essay opens with the writer's grandmother planting a pear tree and ends with the writer picking pears from the same tree. The ending functions mainly to
\ch{frame the essay by returning to its opening image}{concede an opposing view}{introduce a new argument}{provide statistical evidence}
\end{q}

\section{Transitions and signposting}

\textbf{Transitions} signal the logical relationship between ideas. On writing questions, decide the relationship \emph{first}, then pick the word.

\begin{compare}{Relationship}{Transitions}{Test}
Addition & moreover, furthermore, also, in addition & The second idea adds to the first in the same direction. \\
Contrast & however, yet, nevertheless, on the other hand, by contrast & The second idea pushes against the first. \\
Concession & admittedly, granted, of course, although & The writer grants a point before answering it. \\
Cause and effect & therefore, thus, as a result, consequently & The second idea follows from the first. \\
Example & for example, for instance, specifically & The second idea is an instance of the first. \\
Emphasis & indeed, in fact & The second idea intensifies the first. \\
Similarity & similarly, likewise & The second idea parallels the first in another case. \\
Sequence & first, then, finally, meanwhile & Order in time or steps. \\
\end{compare}

\textbf{Signposting} goes beyond single words: a sentence that announces the next move (``The cost is only half the problem'') helps the reader follow the reasoning.

\begin{trap}
``However'' vs.\ ``therefore.'' If the second sentence is a \emph{consequence}, a contrast word is wrong no matter how natural it sounds. Read the two sentences with ``and so'' and with ``but'' to hear which relationship is real.
\end{trap}

\begin{q}{4}{3.4 Transitions}{The second sentence is a consequence of the first.}
``The bridge was closed for repairs for six months. \blank\ , the businesses on the north bank lost most of their customers.'' Which transition best fits the blank?
\chh{However}{Similarly}{For instance}{As a result}
\end{q}

\begin{q}{1}{3.4 Transitions}{The second sentence pushes against the first: contrast.}
``Most residents support the new park. \blank\ , few have volunteered to help maintain it.'' Which transition best fits the blank?
\chh{Nevertheless}{Likewise}{Therefore}{For example}
\end{q}

\section{Reasoning and organization in writing questions}

\begin{terms}
Placement & ``Where should sentence X go?'' Find the sentence it depends on (a pronoun, a ``this,'' an example of a claim) and place it right after. \\
Paragraph order & The paragraph that introduces a term must come before the paragraph that uses it. \\
Topic sentences & The best topic sentence states the paragraph's claim and connects to the one before. \\
Adding a sentence & Keep it only if it advances the reasoning at that point. \\
Concluding & The best conclusion fits the stated goal (call to action, return to the opening, broader implication). \\
\end{terms}

\begin{q}{3}{3.5 Writing: placement}{A sentence beginning ``These delays'' must follow the sentence that describes the delays.}
A writer wants to add the sentence ``These delays cost the average commuter forty hours a year.'' It should be placed
\ch{in the conclusion, after the call to action}{after a sentence about the history of the subway}{immediately after the sentence describing the delays}{at the very beginning of the essay}
\end{q}

\section{Writing the synthesis essay (Q1)}

\subsection{The task}
Q1 presents a topic, a short introduction, and six sources (usually labeled A--F), at least one of them visual, such as a chart, graph, or image. You must develop \emph{your own} position and support it with evidence from \emph{at least three} sources, cited by letter or by the description in parentheses, e.g.\ (Source B) or (Nakamura).

\subsection{During the reading period}
\begin{enumerate}[leftmargin=1.5em]
\item Read the prompt and decide what kind of position it asks for (whether, what factors, to what extent).
\item For each source, write one line: its claim, and whether it supports, complicates, or opposes the position you are forming.
\item Group sources into two or three \emph{ideas}. Your body paragraphs are ideas, not sources.
\end{enumerate}

\subsection{Synthesis, not summary}
Weak synthesis marches through sources one per paragraph (``Source A says\ldots Source B says\ldots''). Strong synthesis puts sources in conversation: two sources support one claim from different angles; a third complicates it; your commentary explains what the combination shows.

\begin{study}{Putting sources in conversation}
``The city's own survey shows that only 12\% of residents use the bus (Source C), but the survey's timing, a summer month when schools were closed, hides the riders who depend on it most: Source E's route data shows student ridership doubles in September. Low average use, in other words, is not low need.''
\end{study}

\subsection{Reaching Row C}
Use a source that \emph{opposes} you and answer it; notice a limitation in a source (its date, its sample, its interests); or weigh competing values the sources represent.

\begin{q}{1}{3.6 Synthesis essay}{Body paragraphs organized by idea, with sources in conversation.}
Which plan is best for a synthesis essay?
\ch{Two or three paragraphs, each built around a claim and drawing on several sources that support or complicate it}{One paragraph per source, beginning with the source that agrees most}{A personal narrative with no sources}{Six paragraphs, one summarizing each source in order}
\end{q}

\begin{q}{2}{3.6 Synthesis essay}{Answering a source that complicates your view shows complexity.}
A writer arguing for later school start times notes that Source D reports bus costs would rise and responds that the district could stagger elementary and high school routes. This move mainly
\ch{summarizes Source D}{shows complexity by answering a source that complicates the argument}{violates the rule against using opposing sources}{abandons the writer's position}
\end{q}

\makekey
\end{document}
